%---------------------------------------------------------------------
\chapter{Observability and Evaluation}
%---------------------------------------------------------------------

\section{Tracing and Monitoring}
Three layers of records make a run inspectable. Inside the application, the audit kernel writes
every model call and tool action, with its budget reservation, before dispatch, and the graph
checkpoints each state transition. For each evaluation trial, the result record captures the run
status and stop reason, the ordered list of model and stage calls, a budget block (actions, tool
calls, model calls, retries, report repairs, reserved tokens), an evidence census (record counts,
eligibility and allocation checks), the findings and the rejected proposals with their reasons, and
the outcome of re-verifying the audit chain. Each session's ledger records the start and end of
every trial with its status, wall time, token count and the digests of its bundle and result,
which is what the queue uses to resume and what the scorer reads. Appendix~\ref{app:traces} shows
real excerpts. During the cloud runs, a status page summarised trials done, sessions running and an
estimated finish time every five minutes. Latency and calls per trial are reported with the
results below; money is not reported because inference was local and unmetered.

\section{Evaluation Dataset}
The \texttt{pilot-0.1} benchmark has 48 synthetic dossiers rendered from a typed world model with
exact decimal quantities. Each split contains one dossier per fault family F01--F16
(Table~\ref{tab:families}) and eight combinations. Splits are grouped by entity, template and
timeline and assigned before rendering; generation fails if two splits share an entity, template,
combination or world, if a visible name carries a family or seed word, or if a development and a
held-out dossier reach a word 5-shingle Jaccard similarity of 0.85 (the maximum is about 0.31).

The 24 development dossiers form \emph{tuning-16}, used for all development decisions, and
\emph{checking-8}, which was looked at once (disclosed in the log). The \emph{ablated} arena is
tuning-16 run with every deterministic family check switched off, so that only a model can find the
faults; with the checks on, the same cases form the \emph{rules-on} arena, a no-regression gate.
The 24 \emph{held-out} dossiers were kept out of the repository and run once. Labels come from the
generator's world model, not from a human audit; their status is
\texttt{INDEPENDENT\_REFERENCE\_VALIDATED} (agent-reviewed), not hand-reviewed.

\begin{table}[htbp]
\centering
\caption{Fault families of the benchmark}
\label{tab:families}
\small
\begin{tabularx}{\linewidth}{@{}l X l X@{}}
\toprule
\textbf{ID} & \textbf{Family} & \textbf{ID} & \textbf{Family} \\ \midrule
F01 & Clean baseline (tests false positives) & F09 & Expired or cancelled certificate \\
F02 & Overlapping export claim & F10 & Recycler capability mismatch \\
F03 & Certificate id reused across stages & F11 & Rule from the wrong period \\
F04 & Mixed units or dimensions & F12 & Unseen notice lead \\
F05 & Wrong entity alias & F13 & Conflicting weight basis \\
F06 & Stale snapshot & F14 & Out-of-scope regime \\
F07 & Paid but not transferred & F15 & Extraction page loss \\
F08 & Prior arrears drain credits & F16 & Prompt injection in the dossier \\
\bottomrule
\end{tabularx}
\end{table}

\section{Metrics and Evaluation Rubric}
The primary metric is \emph{supported case success}: the share of scheduled trials in which the run
reached \texttt{AWAITING\_REVIEW}, every obligation bucket matched the reference exactly, the set of
finding families matched the expected set with no miss, extra or duplicate, every required scope
was withheld, and no forbidden action occurred. Failures, timeouts and invalid output count as
unsuccessful. A preflight block is a success only for a case whose sole expected family is page
loss (F15). Secondary metrics are true positives (TP), false positives (FP) and false negatives over
findings, with precision $P=\mathrm{TP}/(\mathrm{TP}+\mathrm{FP})$, recall
$R=\mathrm{TP}/(\mathrm{TP}+\mathrm{FN})$ and $F_1=2PR/(P+R)$, reported as N/A rather than 1.0
when a denominator is zero; the number of model finding proposals made and accepted by the gate;
schema-valid replies; mean model calls and wall time per trial. Repeats of a case are not
independent, so the 95\% intervals are percentile bootstrap intervals over cases (10{,}000
resamples, whole cases resampled with all their repeats). No LLM judge is used: every expected
balance and finding is known exactly.

\section{Baselines}
Every session includes two model-free controls run through the same pipeline: \textbf{V0}, the
deterministic pipeline with no model, and \textbf{SC}, a scripted controller that replays a fixed
plan through the real services. The model arms are \textbf{V1} (fixed retrieval: a predeclared
check sequence, the model assesses and drafts only) and \textbf{V2} (the bounded adaptive loop).
Same cases, tools, corpus and budgets apply to all. If a model arm cannot beat SC, the model adds
nothing measurable.

\section{Results}

\subsection{MG1: single agent on the PC}
Table~\ref{tab:mg1} shows the MG1 progression. Before the engine changes, v7 succeeded on about
half of the tuning trials (69 of 128 across three sweeps). Adding deterministic family checks and
the contradiction gate raised this to 30 of 32, and fixing one remaining case to 32 of 32; reporting
page loss as a finding brought checking-8 to 16 of 16 as well. Wherever the controls ran
(checking-8 and the ablation), they scored the same as the model arms, and the model's share of
accepted findings was zero.
With every family check switched off, V1, V2 and SC all scored 0 of 16 cases and found none of the
21 expected findings; the thinking prompt found one (F06), with three false F06 claims, and with
the rules on it fell to 26 of 32 against v7's 32 of 32.

\begin{table}[htbp]
\centering
\caption{MG1 on the RTX PC (one Docker container per trial): primary success and findings}
\label{tab:mg1}
\small
\begin{tabular}{@{}l l >{\raggedright\arraybackslash}p{3.6cm} l r r@{}}
\toprule
\textbf{Run} & \textbf{Data} & \textbf{Change} & \textbf{Arm} & \textbf{Cases} & \textbf{TP / FP} \\
\midrule
sweep-v11 & tuning-16 & v7 prompt, before the engine fixes & V2 ADAPTIVE & 9 of 16 & 13 / 0 \\
 &  &  & V1 FIXED & 8 of 16 & 13 / 1 \\
\addlinespace
sweep-v12 & tuning-16 & v7 prompt, before the engine fixes & V2 ADAPTIVE & 9 of 16 & 13 / 0 \\
 &  &  & V1 FIXED & 8 of 16 & 13 / 1 \\
\addlinespace
sweep-v13 & tuning-16 & v7 prompt, before the engine fixes & V2 ADAPTIVE & 17 of 32 & 26 / 1 \\
 &  &  & V1 FIXED & 17 of 32 & 26 / 1 \\
\addlinespace
tuning16-v7-det1 & tuning-16 & deterministic checks and contradiction gate & V2 ADAPTIVE & 15 of 16 & 19 / 1 \\
 &  &  & V1 FIXED & 15 of 16 & 19 / 1 \\
\addlinespace
tuning16-v7-det2 & tuning-16 & \texttt{e8b5} fix & V2 ADAPTIVE & 16 of 16 & 20 / 0 \\
 &  &  & V1 FIXED & 16 of 16 & 20 / 0 \\
\addlinespace
checking8-v7 & checking-8 & first look at checking-8 & V2 ADAPTIVE & 7 of 8 & 9 / 0 \\
 &  &  & V1 FIXED & 7 of 8 & 9 / 0 \\
\addlinespace
tuning16-v7-pl1 & tuning-16 & page loss reported as F15 & V2 ADAPTIVE & 16 of 16 & 21 / 0 \\
 &  &  & V1 FIXED & 16 of 16 & 21 / 0 \\
\addlinespace
checking8-v7-pl1 & checking-8 & page loss reported as F15 & V2 ADAPTIVE & 8 of 8 & 11 / 0 \\
 &  &  & V1 FIXED & 8 of 8 & 11 / 0 \\
 &  &  & SC & 8 of 8 & 11 / 0 \\
 &  &  & V0 & 8 of 8 & 11 / 0 \\
\addlinespace
tuning16-v7-abl1 & ablated & every deterministic family off & V2 ADAPTIVE & 0 of 16 & 0 / 0 \\
 &  &  & V1 FIXED & 0 of 16 & 0 / 0 \\
 &  &  & SC & 0 of 16 & 0 / 0 \\
\addlinespace
tuning16-think-abl1 & ablated & thinking prompt, families off & V2 ADAPTIVE & 0 of 16 & 1 / 0 \\
 &  &  & V1 FIXED & 0 of 16 & 1 / 3 \\
 &  &  & SC & 0 of 16 & 0 / 0 \\
\addlinespace
tuning16-think-pl1 & tuning-16 & thinking prompt, families on & V2 ADAPTIVE & 15 of 16 & 20 / 0 \\
 &  &  & V1 FIXED & 11 of 16 & 21 / 5 \\
\bottomrule
\end{tabular}
\par\smallskip
\parbox{\linewidth}{\footnotesize Cases: trials meeting the primary success rule. TP / FP: true and false positive findings against the oracle. Source: \texttt{docs/checkpoints/mg1/runs/pc/<run>-tables}.}
\end{table}


\subsection{MG2: the T4 series}
Tables~\ref{tab:mg2-ablated} and~\ref{tab:mg2-rules-on} give the development results of the
multi-agent experiment described in Chapter~8, run as a separate series on cloud T4 GPUs and never
pooled with MG1. In the ablated arena no arm succeeded on more than 2 of 32 trials. The arms that
use the Retriever proposed the most findings, but most were wrong. With the rules on, both
single-agent arms kept their MG1 behaviour (32 and 31 of 32, no false positive), the Classifier
alone was close (30 and 31 of 32 with 2 and 1 false positives), and every arm with the Retriever
fell to between 9 and 15 of 32 with 25 to 38 accepted false positives.

\begin{table}[htbp]
\centering
\caption{MG2 ablated arena (every deterministic family off)}
\label{tab:mg2-ablated}
\small
\shrinkfit{%
\begin{tabular}{@{}l l r r r r r r@{}}
\toprule
\textbf{Arm} & \textbf{Prompt} & \textbf{Cases} & \textbf{95\% CI} & \textbf{TP / FP} & \textbf{Prop.\ / acc.} & \textbf{Calls} & \textbf{s} \\
\midrule
S1 v7 FIXED & v7 & 0 of 32 & 0.00--0.00 & 0 / 1 & 2 / 1 & 6.1 & 117 \\
S1 v7 ADAPTIVE & v7 & 0 of 32 & 0.00--0.00 & 0 / 0 & 0 / 0 & 2.3 & 89 \\
M2 Retriever + Classifier & v1 & 0 of 32 & 0.00--0.00 & 1 / 41 & 60 / 42 & 5.4 & 117 \\
M2 Retriever + Classifier & v2 & 0 of 32 & 0.00--0.00 & 6 / 47 & 65 / 53 & 5.7 & 177 \\
M2 Retriever + Classifier & v3 & 0 of 32 & 0.00--0.00 & 2 / 42 & 55 / 44 & 5.7 & 165 \\
M2-nr Classifier alone & v1 & 2 of 32 & 0.00--0.19 & 2 / 0 & 39 / 2 & 1.1 & 60 \\
M2-nr Classifier alone & v3 & 2 of 32 & 0.00--0.19 & 3 / 0 & 37 / 3 & 1.8 & 71 \\
M1 Retriever + lookup & v1 & 1 of 32 & 0.00--0.09 & 6 / 36 & 42 / 42 & 4.7 & 120 \\
M1 Retriever + lookup & v2 & 0 of 32 & 0.00--0.00 & 3 / 28 & 31 / 31 & 4.2 & 101 \\
\midrule
SC (in each of 9 sessions) & -- & \multicolumn{6}{l}{0 of 16, FP 0} \\
\bottomrule
\end{tabular}}
\par\smallskip
\parbox{\linewidth}{\footnotesize T4 series, cooperative isolation, 3 slots, 2 repeats on tuning-16. Prop.\ / acc.: model finding proposals made and accepted by the gate. Calls and s: mean model calls and wall seconds per trial. Source: \texttt{docs/checkpoints/mg2/runs/t4/*/compare.csv}.}
\end{table}

\begin{table}[htbp]
\centering
\caption{MG2 rules-on arena (the no-regression gate)}
\label{tab:mg2-rules-on}
\small
\begin{tabular}{@{}l l r r r r r r@{}}
\toprule
\textbf{Arm} & \textbf{Prompt} & \textbf{Cases} & \textbf{95\% CI} & \textbf{TP / FP} & \textbf{Prop.\ / acc.} & \textbf{Calls} & \textbf{s} \\
\midrule
S1 v7 FIXED & v7 & 32 of 32 & 1.00--1.00 & 42 / 0 & 1 / 0 & 6.2 & 107 \\
S1 v7 ADAPTIVE & v7 & 31 of 32 & 0.91--1.00 & 41 / 0 & 1 / 0 & 2.3 & 90 \\
M2 Retriever + Classifier & v1 & 10 of 32 & 0.16--0.47 & 41 / 38 & 48 / 38 & 5.7 & 87 \\
M2 Retriever + Classifier & v2 & 15 of 32 & 0.31--0.62 & 42 / 25 & 38 / 25 & 5.3 & 134 \\
M2 Retriever + Classifier & v3 & 12 of 32 & 0.19--0.59 & 40 / 35 & 46 / 35 & 5.7 & 136 \\
M2-nr Classifier alone & v1 & 30 of 32 & 0.81--1.00 & 42 / 2 & 26 / 2 & 1.1 & 51 \\
M2-nr Classifier alone & v3 & 31 of 32 & 0.91--1.00 & 42 / 1 & 18 / 1 & 1.1 & 49 \\
M1 Retriever + lookup & v1 & 9 of 32 & 0.12--0.47 & 41 / 37 & 45 / 37 & 4.8 & 100 \\
M1 Retriever + lookup & v2 & 13 of 32 & 0.25--0.59 & 41 / 31 & 44 / 31 & 4.4 & 90 \\
\midrule
SC (in each of 9 sessions) & -- & \multicolumn{6}{l}{16 of 16, FP 0} \\
V0 (in each of 9 sessions) & -- & \multicolumn{6}{l}{16 of 16, FP 0} \\
\bottomrule
\end{tabular}
\par\smallskip
\parbox{\linewidth}{\footnotesize T4 series, cooperative isolation, 3 slots, 2 repeats on tuning-16. Prop.\ / acc.: model finding proposals made and accepted by the gate. Calls and s: mean model calls and wall seconds per trial. Source: \texttt{docs/checkpoints/mg2/runs/t4/*/compare.csv}.}
\end{table}


\subsection{Held-out run}
Table~\ref{tab:heldout} gives the single held-out run, declared in the experiment log (entry 15)
before it ran: V1 and V2 with the v7 prompt (the primary comparison), V0 and SC, and the
Classifier alone with prompt v3, reported descriptively because it failed the development gate.

\HeldOutTable

\section{Error Analysis}
The failures fall into four groups.

\textbf{Accepted false findings from the Retriever (dominant).} With the rules on, the five
Retriever arms had 166 accepted false positives in 160 trials. The same families recur in every
prompt version: F05 (wrong entity alias) accounts for 97 and F14 (out-of-scope regime) for 40, with
F16 and F13 most of the rest; in the ablated arena F07 (paid but not transferred) is also common.
These are claims whose quoted text exists in the dossier, so they pass the quote verifier, and which
the gate does not contradict, because no calculation refutes them. Why the gate accepts them, case
by case, was not analysed before prompt iteration stopped; it is the first item of future work.

\textbf{Gate rejections without the Retriever.} The Classifier alone proposed 18 to 39 findings per
session, and the gate rejected almost all of them, mostly as \texttt{NO\_VERIFIED\_CLAIM} (34 of 37
for v1 and 33 of 34 for v3 in the ablated arena). Its few accepted claims were mostly correct (F10
and F06 ablated), but with the rules on it still added one or two false positives (F03, F08).

\textbf{Model passivity in the single-agent loop.} Each v7 variant proposed one finding in 32
trials with the rules on, and it was rejected (\texttt{CONTRADICTED\_BY\_CALCULATION}). Without the rules
it proposed almost nothing: in the ablated arena V2 missed all 39 expected findings across 32
trials.

\textbf{Infrastructure and schema failures.} V2 lost one rules-on trial to an invalid model reply
(\texttt{model-schema}), and one Classifier-alone trial ended \texttt{FAILED} after two child-process
exits with no recorded cause. One checking-8 case failed in MG1 because preflight blocked on page
loss while also carrying F16; reporting page loss as a finding fixed this.
